%%% A Closed-Loop Theorem for SGD's Noise-Driven Selection
%%% in the Non-Interpolation Regime
%%%
%%% Type A paper: theorem with rigorous proof.
%%% Companion to paper-a (interpolation regime, discrete Lyapunov)
%%% and paper-b (failure modes F1-F4).

\documentclass[11pt,a4paper]{article}

%% ---------- packages ----------
\usepackage[utf8]{inputenc}
\usepackage[T1]{fontenc}
\usepackage{amsmath,amssymb,amsthm,mathtools}
\usepackage[margin=2.5cm]{geometry}
\usepackage{hyperref}
\usepackage{cleveref}
\usepackage{booktabs}
\usepackage{enumitem}

%% ---------- theorem environments ----------
\newtheorem{theorem}{Theorem}[section]
\newtheorem{proposition}[theorem]{Proposition}
\newtheorem{lemma}[theorem]{Lemma}
\newtheorem{corollary}[theorem]{Corollary}
\newtheorem{conjecture}[theorem]{Conjecture}
\theoremstyle{definition}
\newtheorem{definition}[theorem]{Definition}
\newtheorem{assumption}{Assumption}
\newtheorem{example}[theorem]{Example}
\theoremstyle{remark}
\newtheorem{remark}[theorem]{Remark}

%% ---------- macros ----------
\DeclareMathOperator{\tr}{tr}
\DeclareMathOperator{\spn}{span}
\DeclareMathOperator{\diag}{diag}
\DeclareMathOperator{\curl}{curl}
\DeclareMathOperator{\Var}{Var}
\DeclareMathOperator{\Cov}{Cov}
\newcommand{\R}{\mathbb{R}}
\newcommand{\E}{\mathbb{E}}
\newcommand{\Prob}{\mathbb{P}}
\newcommand{\dif}{\,\mathrm{d}}
\newcommand{\Veff}{V_{\mathrm{eff}}}
\newcommand{\Gap}{\mathrm{Gap}}
\newcommand{\LOO}{\mathrm{LOO}}
\newcommand{\norm}[1]{\left\|#1\right\|}
\newcommand{\inner}[2]{\langle #1,#2\rangle}

%% ---------- metadata ----------
\title{\Large A Closed-Loop Theorem for SGD's Noise-Driven Selection\\[4pt]
       \large in the Non-Interpolation Regime}
\author{Lightman Chang\\
        Independent Researcher\\
        \texttt{lightman.chang@gmail.com}}
\date{}

%% =================================================================
\begin{document}
\maketitle

%% =================================================================
%% ABSTRACT
%% =================================================================
\begin{abstract}
We study the implicit regularization of stochastic gradient descent (SGD) in
the \emph{non-interpolation} regime, where the gradient noise covariance
$\Sigma(\theta^*)$ at a local minimum is strictly positive and the
continuous-time Fokker--Planck approximation is non-degenerate.  Our main
result (\Cref{thm:closed-loop}) is a closed-loop theorem with three
quantitatively linked components: (i) the SGD stationary measure on a
multi-well potential admits an explicit Gibbs-type formula whose
relative weights at two minima are governed by the ratios
$\Sigma_A/\Sigma_B$ and $H_A/H_B$; (ii) the leave-one-out generalization gap
at a quadratic minimum equals $\Sigma(\theta^*)/(H(\theta^*)(n-1))$ up to
$O(n^{-2})$; (iii) consequently $\Sigma_A/H_A < \Sigma_B/H_B$ implies that
SGD concentrates on the basin of $\theta_A^*$ \emph{and} that $\theta_A^*$
has a strictly smaller test-loss bound than $\theta_B^*$, with explicit
margin.  We extend the dynamic part to multivariate parameter spaces via the
Freidlin--Wentzell quasipotential, working out a two-dimensional example with
state-dependent diffusion in which detailed balance fails but the
quasipotential admits a closed-form Hamilton--Jacobi expansion
(\Cref{thm:multi}).  We prove that the ratio $\Sigma/H$, and its
multivariate generalization $\tr(H^{-1}\Sigma)$, are invariant under smooth
diffeomorphisms of the parameter space (\Cref{prop:reparam}), which makes
the selection criterion intrinsic.  The result is the non-interpolation
counterpart of the discrete Lyapunov closed loop established in
\cite{paperA} for the interpolation regime, and its sharpness is exactly
delimited by the four failure modes catalogued in \cite{paperB}; in
particular, the implication (iii) is reversed by failure mode~F4, and the
interplay is made precise.
\end{abstract}

\noindent\textbf{MSC 2020:} 60H10 (primary); 62L20, 60F10, 68T07,
49L25, 60J60, 65C05.

%% =================================================================
\section{Introduction}\label{sec:intro}
%% =================================================================

\paragraph{The problem.}
Stochastic gradient descent (SGD) is the workhorse optimizer of modern
machine learning, and a long line of work has tried to explain why it
selects, among the many local minima of a non-convex training loss, those
that generalize well to unseen data.  In the \emph{non-interpolation
regime}---where the training loss does not vanish, so the per-sample
gradients $\nabla\ell_i(\theta^*)$ at a local minimum~$\theta^*$ are not all
zero---a popular explanation models SGD by the It\^o stochastic
differential equation (SDE)
\[
   d\theta_t \;=\; -\nabla L(\theta_t)\dif t
                  \;+\; \sqrt{\tau\,\Sigma(\theta_t)}\,dW_t,
   \qquad
   \tau \;=\; \eta/B,
\]
where $\eta$ is the learning rate, $B$ is the batch size and
$\Sigma(\theta)$ is the per-sample gradient covariance
\cite{LiTaiE2017,SmithLe2018,MandtHoffmanBlei2017,ChaudhariSoatto2018}.  The
multiplicative structure of $\Sigma$ then biases the Fokker--Planck
stationary measure toward minima of low noise covariance.  This is
appealing as a heuristic, but two questions have remained open at a fully
rigorous level:
\begin{enumerate}[label=(Q\arabic*),leftmargin=2.6em]
\item \label{Q1}
   Can the dynamic preference of SGD for a particular local minimum be
   quantified by an explicit, parameter-free formula in terms of the
   curvature $H(\theta^*)$ and the noise covariance $\Sigma(\theta^*)$?
\item \label{Q2}
   Does this dynamic preference align with a sample-based notion of
   generalization---e.g.\ leave-one-out (LOO) stability---and if so,
   through which intrinsic geometric quantity?
\end{enumerate}

\paragraph{What we do.}
We answer \ref{Q1} and \ref{Q2} in the affirmative under explicit assumptions
that we state quantitatively.  Specifically, we prove a single closed-loop
theorem that links three components: the Fokker--Planck stationary
distribution (selection), the leave-one-out gap (generalization), and the
implication that connects them.  All three reduce, in the symmetric
two-well one-dimensional case, to the single intrinsic quantity
$\Sigma(\theta^*)/H(\theta^*)$.  The qualitative picture has been suggested
in earlier work \cite{ChaudhariSoatto2018,SmithLe2018,HaoChen2021}; the
contribution of the present paper is to make every step quantitative and
self-contained, including a multivariate extension via the
Freidlin--Wentzell quasipotential and a complete reparameterization
invariance proof.

\paragraph{Delta from prior work.}
The components of \Cref{thm:closed-loop} are not all new in isolation.
\emph{Part~I} (the one-dimensional Gibbs-type stationary density of the
Fokker--Planck equation with multiplicative noise) is textbook material
in stochastic processes \cite{Pavliotis2014} and was applied to SGD
explicitly by Mandt--Hoffman--Blei \cite[\S3]{MandtHoffmanBlei2017} and
others; the closed-form
$p_\tau\propto\Sigma^{-1}\exp(-(2/\tau)U)$ is folklore in this
literature.  \emph{Part~II} (the leave-one-out generalization gap
$\Sigma/(H(n-1))$ at a quadratic minimum) is a classical second-order
expansion of stability \cite{BousquetElisseeff2002} and is implicit in
optimization-theoretic surveys \cite{BottouCurtisNocedal2018}.
\emph{The contribution of the present paper} is to bind these three
ingredients together into a single closed loop: stationary mass,
LOO gap, and reparameterization invariance through the common
intrinsic quantity $\Sigma/H$ (resp.\ $\tr(H^{-1}\Sigma)$ in higher
dimensions).  In particular, we (i) state the loop quantitatively
under explicit assumptions, (ii) make the boundary with each of the
four failure modes F1--F4 precise (\Cref{sec:discussion}), and (iii)
extend the dynamic part to the multivariate setting where detailed
balance fails, via a worked Freidlin--Wentzell quasipotential
(\Cref{thm:multi}).  The \emph{closed-loop bind} (Parts I+II+III+
reparameterization), not its individual components, is the present
paper's contribution.

\paragraph{Position relative to companion papers.}
This paper is the third in a series.  Paper~A \cite{paperA} treats the
\emph{interpolation regime}, where every interpolating solution has
$\Sigma(\theta^*)=0$, the SDE approximation degenerates, and a discrete
Lyapunov-exponent analysis must replace the Fokker--Planck argument.
Paper~B \cite{paperB} catalogues four failure modes (F1: insufficient
mixing; F2: discrete instability; F3: degenerate selection; F4:
noise--generalization reversal) under which the closed loop breaks.  The
present paper-C is the positive counterpart in the non-interpolation regime:
it states and proves the closed loop precisely under the negation of those
four failure modes, and isolates the hypothesis on which the implication
between selection and generalization rests, namely the alignment of low
$\Sigma/H$ with low test loss (the negation of F4).

\paragraph{Main results, informally.}
Let $\theta_A^*$ and $\theta_B^*$ be two locally quadratic minima of $L$
with curvatures $H_A,H_B>0$ and noise covariances $\Sigma_A,\Sigma_B>0$.
Under explicit ergodicity assumptions:
\begin{itemize}
\item (\Cref{thm:closed-loop} Part~I) The SGD stationary density is
   $p(\theta) \propto \Sigma(\theta)^{-1}\exp(-(2/\tau)\,U(\theta))$
   for an explicit drift potential $U$, and in the symmetric one-dimensional
   case $p(\theta_A^*)/p(\theta_B^*) = \Sigma_B/\Sigma_A$.
\item (\Cref{thm:closed-loop} Part~II) The LOO generalization gap at a
   quadratic minimum $\theta^*$ satisfies
   $\Gap_{\LOO}(\theta^*) \;=\; \Sigma(\theta^*)/(H(\theta^*)(n-1)) \;+\;
                                                                  O(n^{-2})$.
\item (\Cref{thm:closed-loop} Part~III) If $\Sigma_A/H_A < \Sigma_B/H_B$,
   then the SGD-stationary measure concentrates on the basin of $\theta_A^*$
   in the sense of \Cref{thm:closed-loop}, and $\theta_A^*$'s test-loss
   bound is strictly smaller than that of $\theta_B^*$ with explicit margin
   $\bigl(\Sigma_B/H_B - \Sigma_A/H_A\bigr)/(n-1)$.
\item (\Cref{thm:multi}) For a state-dependent multivariate diffusion in
   which detailed balance fails, the Freidlin--Wentzell quasipotential
   exists and is the unique viscosity solution of an explicit
   Hamilton--Jacobi equation.  We give a worked $2$D example for which
   the leading correction to the naive potential is computed in closed
   form.
\item (\Cref{prop:reparam}) The ratio $\Sigma/H$, and its multivariate
   generalization $\tr(H^{-1}\Sigma)$, are invariant under smooth
   diffeomorphisms of the parameter space.
\end{itemize}

\paragraph{Organization.}
\Cref{sec:prelim} fixes the notation and lists the standing assumptions.
\Cref{sec:theorem-statement} states the main theorem.
\Cref{sec:proof-1,sec:proof-2,sec:proof-3} contain the three parts of the
proof.  \Cref{sec:multi} contains the multivariate extension.
\Cref{sec:reparam} contains the reparameterization invariance proof.
\Cref{sec:discussion} discusses the boundary with paper~B, the relation to
paper~A, and the open problems.

%% =================================================================
\section{Preliminaries}\label{sec:prelim}
%% =================================================================

\subsection{The empirical risk and SGD}\label{sec:sgd}

Let $S = \{z_1,\dots,z_n\} \subset \mathcal{Z}$ be a training set drawn
i.i.d.\ from a distribution $\mathcal{D}$.  Fix a parameter space
$\Theta \subseteq \R^d$ and a per-sample loss
$\ell:\Theta \times \mathcal{Z}\to \R$ that is $C^2$ in $\theta$.  Write
\[
   L(\theta) \;=\; \frac{1}{n}\sum_{i=1}^n \ell_i(\theta),
   \qquad
   \ell_i(\theta) \;:=\; \ell(\theta,z_i).
\]
SGD is the recursion
\begin{equation}\label{eq:sgd-recursion}
   \theta_{t+1} \;=\; \theta_t \;-\; \eta\,\nabla \ell_{I_t}(\theta_t),
   \qquad I_t \sim \mathrm{Uniform}\{1,\dots,n\},\;\; \text{i.i.d.}
\end{equation}
We work with constant learning rate $\eta>0$.  The mini-batch case
$B\geq 1$ is identical with $\tau:=\eta/B$ in place of $\eta$ throughout.

\subsection{Curvature and gradient noise}

\begin{definition}[Curvature and noise covariance]\label{def:H-Sigma}
   At any $\theta\in\Theta$,
   \begin{align*}
      H(\theta) &\;:=\; \nabla^2 L(\theta) \;=\;
                       \frac{1}{n}\sum_{i=1}^n \nabla^2 \ell_i(\theta),\\
      \Sigma(\theta) &\;:=\;
         \frac{1}{n}\sum_{i=1}^n
            \bigl(\nabla\ell_i(\theta) - \nabla L(\theta)\bigr)
            \bigl(\nabla\ell_i(\theta) - \nabla L(\theta)\bigr)^{\!\top}.
   \end{align*}
   When $d=1$ both reduce to scalars.  At a critical point $\theta^*$
   ($\nabla L(\theta^*)=0$), the noise covariance simplifies to
   $\Sigma(\theta^*) = \frac{1}{n}\sum_{i=1}^n \nabla\ell_i(\theta^*)
                                              \nabla\ell_i(\theta^*)^\top$.
\end{definition}

\begin{definition}[Three-level convergence]\label{def:three-level}
   We distinguish:
   \begin{itemize}[leftmargin=2em]
   \item[\textbf{(C1)}] \emph{Loss convergence}: $L(\theta_t) \to L^*$
                       almost surely, for some $L^*$.
   \item[\textbf{(C2)}] \emph{Distributional convergence}: there is a
                       probability measure $\mu_\eta$ on $\Theta$ such that
                       the law of $\theta_t$ converges weakly to
                       $\mu_\eta$ as $t \to \infty$.
   \item[\textbf{(C3)}] \emph{Absorption convergence}: $\theta_t\to\theta^*$
                       almost surely for some critical point $\theta^*$.
   \end{itemize}
\end{definition}

In the non-interpolation regime, (C3) does not hold for SGD with constant
$\eta>0$: the iterates fluctuate in a basin rather than converge to a
single point.  The natural mode of convergence is (C2), and we shall
characterize the limit $\mu_\eta$.

\subsection{The Fokker--Planck approximation}

For small $\eta$, the recursion \eqref{eq:sgd-recursion} is approximated
\cite{LiTaiE2017,MandtHoffmanBlei2017,SmithLe2018} by the It\^o SDE
\begin{equation}\label{eq:sgd-sde}
   d\theta_t \;=\; -\nabla L(\theta_t)\,\dif t
                  \;+\;\sqrt{\tau\,\Sigma(\theta_t)}\,dW_t,
   \qquad \tau = \eta/B.
\end{equation}
The corresponding Fokker--Planck equation (FPE) for the density $p(t,\theta)$ is
\begin{equation}\label{eq:fpe}
   \partial_t p \;=\;
       \nabla\!\cdot\!\bigl(\nabla L\,p\bigr)
   \;+\; \frac{\tau}{2}\,
       \nabla\!\cdot\!\nabla\!\cdot\!\bigl(\Sigma\,p\bigr).
\end{equation}
We take \eqref{eq:fpe} as the working model for the density $\mu_\eta$ and
state our assumptions accordingly.  The standard estimate that justifies
\eqref{eq:sgd-sde}--\eqref{eq:fpe} is that any time-marginal
$\E[\phi(\theta_T)]$ for $\phi\in C^4_b$ agrees up to $O(\eta^2)$ with the
SDE prediction at horizon $T = N\eta$ uniformly for bounded $N$
\cite{LiTaiE2017}.  Throughout this paper we therefore identify SGD's
long-time behaviour with the stationary measure of \eqref{eq:fpe}.

\subsection{Standing assumptions}

\begin{assumption}[Smoothness]\label{A:C2}
   $L\in C^3(\Theta)$ and $\ell_i\in C^3(\Theta)$ for $i=1,\dots,n$.
   The $C^3$ regularity of each $\ell_i$ is required by the Taylor remainder
   $O(\norm{\theta^*_{-j}-\theta^*}^3)$ used in the proof of Part~II
   (\Cref{lem:pointwise-gap}).
\end{assumption}

\begin{assumption}[Coercivity]\label{A:coerc}
   $L(\theta)\to+\infty$ as $\norm{\theta}\to\infty$, with growth at most
   polynomial: there is $r>0$ such that $\norm{\nabla L(\theta)}
   \leq C(1+\norm{\theta}^r)$.
\end{assumption}

\begin{assumption}[Uniform ellipticity of noise]\label{A:ellip}
   There are constants $0 < \sigma_{\min}^2 \leq \sigma_{\max}^2 < \infty$
   such that for every $\theta$ and every unit vector $v\in\R^d$,
   \[
      \sigma_{\min}^2 \;\leq\; v^{\!\top}\Sigma(\theta)\,v
                    \;\leq\; \sigma_{\max}^2.
   \]
\end{assumption}

\begin{assumption}[Two locally quadratic minima]\label{A:two-wells}
   $L$ has exactly two local minima $\theta_A^*,\theta_B^*$ with strict
   positive-definite Hessians $H_A:=H(\theta_A^*)\succ 0$ and
   $H_B:=H(\theta_B^*)\succ 0$, and exactly one saddle $\theta_S$ on the
   minimum-energy path between them.  We write
   $\Sigma_A:=\Sigma(\theta_A^*)$, $\Sigma_B:=\Sigma(\theta_B^*)$.
\end{assumption}

\begin{assumption}[Selection hypothesis]\label{A:sel}
   $\Sigma_A/H_A < \Sigma_B/H_B$ (in $d=1$ as scalars; in higher dimensions
   with the trace generalization of \Cref{prop:reparam}).
\end{assumption}

\begin{assumption}[F4-negation: LOO aligns with population]\label{A:F4-neg}
   The LOO gap is an unbiased proxy for the population--training gap with
   no model-misspecification reversal: at both $\theta_A^*$ and $\theta_B^*$,
   $\E[L_{\mathrm{pop}}(\theta^*) - L(\theta^*)] = \E[\Gap_{\LOO}(\theta^*)]
                                                          + o(1/n)$,
   and the sign of the difference of population gaps between the two
   minima agrees with the sign of the difference of LOO gaps.
   Equivalently, the failure mode F4 of \cite{paperB} is inactive: low
   $\Sigma/H$ corresponds to low population loss rather than memorization
   under data--architecture misspecification.
\end{assumption}

\Cref{A:C2}--\Cref{A:ellip} are mild regularity conditions standard in
diffusion theory \cite{Pavliotis2014}.  Coercivity (\Cref{A:coerc}) and
uniform ellipticity (\Cref{A:ellip}) together imply that the FPE
\eqref{eq:fpe} admits a unique invariant probability density.
\Cref{A:two-wells} restricts attention to the simplest non-trivial
landscape.  \Cref{A:sel} is the hypothesis on which the third part of the
main theorem rests; it is the negation of failure mode~F3
of \cite{paperB} (when the inequality is strict and the alignment with
generalization is preserved, F4 is also negated).

\subsection{Leave-one-out perturbation}

\begin{definition}[LOO empirical risk and minimizer]\label{def:LOO}
   For $j\in\{1,\dots,n\}$ define the leave-one-out empirical risk
   \[
      L_{-j}(\theta) \;:=\; \frac{1}{n-1}\sum_{i\neq j} \ell_i(\theta),
   \]
   and let $\theta^*_{-j}$ be its local minimizer in the same basin as
   $\theta^*$.  The LOO generalization gap at $\theta^*$ is
   \[
      \Gap_{\LOO}(\theta^*) \;:=\;
      \frac{1}{n}\sum_{j=1}^n
      \bigl(\ell_j(\theta^*_{-j}) - \ell_j(\theta^*)\bigr).
   \]
\end{definition}

A standard consequence of $\E[\Gap_{\LOO}] \approx \E[L_{\text{pop}}(\theta^*)
   - L(\theta^*)]$ \cite{BousquetElisseeff2002,MohriRostamizadehTalwalkar2018}
makes $\Gap_{\LOO}$ a sample-based proxy for the population gap.

%% =================================================================
\section{The closed-loop theorem}\label{sec:theorem-statement}
%% =================================================================

\begin{theorem}[Non-interpolation closed loop]\label{thm:closed-loop}
   Adopt \Cref{A:C2}--\Cref{A:two-wells}, set $\tau = \eta/B$, and consider
   the It\^o SDE \eqref{eq:sgd-sde}.  Then:
   \begin{enumerate}[label=\textnormal{(\Roman*)},leftmargin=2.6em]
   \item \emph{(Stationary measure.)}  The FPE \eqref{eq:fpe} admits a
       unique probability invariant density $p_\tau(\theta)$.  In dimension
       $d=1$,
       \begin{equation}\label{eq:stationary-1d}
          p_\tau(\theta) \;=\; \frac{1}{Z_\tau\,\Sigma(\theta)}\,
              \exp\!\Bigl(-\frac{2}{\tau}\,U(\theta)\Bigr),
          \quad
          U(\theta) \;:=\; \int_{\theta_0}^\theta
                 \frac{L'(s)}{\Sigma(s)}\,ds,
       \end{equation}
       where $\theta_0$ is any reference point and $Z_\tau$ normalizes
       $p_\tau$.  Let $B_A,B_B\subset\R$ denote the open basins of
       attraction of $\theta_A^*,\theta_B^*$ for $-L'$ (separated by
       the saddle $\theta_S$), and write the basin masses as
       $M_\tau(A) := \int_{B_A} p_\tau(\theta)\,d\theta$,
       $M_\tau(B) := \int_{B_B} p_\tau(\theta)\,d\theta$.  Applying
       Laplace's method to \eqref{eq:stationary-1d} as $\tau\to 0$ gives
       \begin{equation}\label{eq:ratio-laplace}
          \frac{M_\tau(A)}{M_\tau(B)}
          \;=\; \sqrt{\frac{H_B}{H_A}}\,\frac{\Sigma_B}{\Sigma_A}\,
                \exp\!\Bigl(-\tfrac{2}{\tau}\bigl(U(\theta_A^*)
                                                -U(\theta_B^*)\bigr)\Bigr)
                \,(1 + O(\tau)).
       \end{equation}
       In the symmetric two-well case ($L$ and $\Sigma$ symmetric about
       $(\theta_A^*+\theta_B^*)/2$, so that $H_A=H_B$ and
       $U(\theta_A^*)=U(\theta_B^*)$),
       \begin{equation}\label{eq:ratio-1d}
          \frac{M_\tau(A)}{M_\tau(B)}
          \;=\; \frac{\Sigma_B}{\Sigma_A}\,(1+O(\tau))
          \;\;\bigl(>1 \text{ iff } \Sigma_A<\Sigma_B\bigr).
       \end{equation}
       In the general (asymmetric) one-dimensional case the basin-mass
       ratio is given by \eqref{eq:ratio-laplace}; the corresponding
       \emph{pointwise} density ratio is
       \begin{equation}\label{eq:ratio-asym}
          \frac{p_\tau(\theta_A^*)}{p_\tau(\theta_B^*)}
          \;=\; \frac{\Sigma_B}{\Sigma_A}\,
                \exp\!\Bigl(-\tfrac{2}{\tau}\bigl(U(\theta_A^*)
                                                -U(\theta_B^*)\bigr)\Bigr).
       \end{equation}
   \item \emph{(LOO gap.)}  At a locally quadratic minimum
       $\theta^*$ with $H(\theta^*)>0$, under \Cref{A:C2},
       \begin{equation}\label{eq:LOO-gap}
          \Gap_{\LOO}(\theta^*) \;=\;
              \frac{\Sigma(\theta^*)}{(n-1)\,H(\theta^*)}
              \;+\; O\!\Bigl(\frac{1}{n^2}\Bigr),
       \end{equation}
       where the constant in $O(n^{-2})$ depends only on
       $\sup_{j}\bigl(\norm{\nabla \ell_j(\theta^*)}^4 +
       \norm{\nabla^3 \ell_j(\theta^*)}\bigr)$ and $H(\theta^*)^{-1}$.
   \item \emph{(Implication.)}  Assume in addition \Cref{A:sel}, i.e.\
       $\Sigma_A/H_A < \Sigma_B/H_B$, and the equal-loss two-well setting
       $L(\theta_A^*) = L(\theta_B^*)$ so that the basin-mass formula
       \eqref{eq:ratio-laplace} applies.  Then
       \begin{enumerate}[label=\textnormal{(\roman*)}]
       \item the stationary mass concentrates on the basin of $\theta_A^*$:
           for $\tau$ sufficiently small,
           $M_\tau(A)/M_\tau(B) > 1$;
       \item the LOO gap at well $A$ is strictly smaller than at well $B$:
           \begin{equation}\label{eq:gap-margin}
              \Gap_{\LOO}(\theta_B^*) - \Gap_{\LOO}(\theta_A^*)
              \;=\; \frac{1}{n-1}\Bigl(
                       \frac{\Sigma_B}{H_B} - \frac{\Sigma_A}{H_A}\Bigr)
                  \;+\; O(n^{-2}) \;>\; 0;
           \end{equation}
       \item under the additional assumption \Cref{A:F4-neg}
           (LOO is an unbiased proxy for the population gap, no F4
           reversal), the test-loss bound at~$A$ is below that at~$B$
           by the same margin (up to $O(n^{-2})$).
       \end{enumerate}
   \end{enumerate}
\end{theorem}

The remainder of \Cref{sec:proof-1,sec:proof-2,sec:proof-3} contains the
proof of \Cref{thm:closed-loop} in three parts.

%% =================================================================
\section{Proof of \Cref{thm:closed-loop}: Part~I (stationary measure)}
\label{sec:proof-1}
%% =================================================================

We work in dimension $d=1$ throughout this section; the multivariate
extension is treated separately in \Cref{sec:multi} via the
Freidlin--Wentzell quasipotential.  Write $L,\Sigma:\R\to\R$ and
$L'(\theta) = dL/d\theta$.

\subsection{Existence and uniqueness of the stationary density}

\begin{lemma}[Existence and uniqueness]\label{lem:exist-unique}
   Under \Cref{A:C2}--\Cref{A:ellip}, the FPE \eqref{eq:fpe} has a unique
   invariant probability density $p_\tau\in C^2(\R)$.
\end{lemma}

\begin{proof}
   The drift $-L'$ is continuous; by \Cref{A:coerc} there are constants
   $a,b>0$ such that $-\theta\,L'(\theta) \leq -a\theta^2 + b$ outside a
   compact set, which is a Lyapunov condition for the SDE
   \eqref{eq:sgd-sde}.  Combined with the uniform ellipticity
   \Cref{A:ellip}, classical results (\cite[Theorem~6.16]{Pavliotis2014})
   give the existence of a unique invariant probability measure with a
   $C^2$ density $p_\tau$, which solves the stationary FPE
   $\partial_t p_\tau = 0$ in \eqref{eq:fpe}.
\end{proof}

\subsection{Closed form of the one-dimensional stationary density}

\begin{lemma}[Stationary density in $d=1$]\label{lem:stationary-1d}
   In dimension $d=1$, the stationary density satisfies
   \begin{equation}\label{eq:stationary-1d-restate}
      p_\tau(\theta)
      \;=\; \frac{1}{Z_\tau\,\Sigma(\theta)}\,
            \exp\!\Bigl(-\frac{2}{\tau}\,U(\theta)\Bigr),
      \qquad U(\theta) \;=\; \int_{\theta_0}^\theta
                                       \frac{L'(s)}{\Sigma(s)}\,ds.
   \end{equation}
   Equivalently, $p_\tau(\theta) \propto \exp(-(2/\tau)\Veff(\theta))$ with
   effective potential
   \begin{equation}\label{eq:Veff}
      \Veff(\theta) \;:=\; U(\theta) \;+\; \frac{\tau}{2}\,\ln \Sigma(\theta).
   \end{equation}
\end{lemma}

\begin{proof}
   The stationary FPE in $d=1$ reads
   \[
      0 \;=\; \frac{d}{d\theta}\Bigl[L'(\theta)\,p(\theta)\Bigr]
            + \frac{\tau}{2}\,
              \frac{d^2}{d\theta^2}\Bigl[\Sigma(\theta)\,p(\theta)\Bigr].
   \]
   Define the probability current
   $J(\theta) := -L'(\theta)\,p(\theta) - \frac{\tau}{2}\,
                    \frac{d}{d\theta}\bigl[\Sigma(\theta)\,p(\theta)\bigr]$.
   Then the FPE is $dJ/d\theta = 0$, so $J$ is a constant.  By
   \Cref{A:coerc} and \Cref{A:ellip}, $p_\tau \in L^1(\R)$ with
   exponential tails, hence $p_\tau(\theta)\to 0$ as $\norm\theta\to\infty$
   and the same holds for $\Sigma\,p_\tau$ and its derivatives.
   Sending $\theta\to\infty$ along a sequence on which the limit exists
   forces $J\equiv 0$.

   Setting $J\equiv 0$:
   \[
      L'(\theta)\,p(\theta)
      \;+\; \frac{\tau}{2}\,\Sigma'(\theta)\,p(\theta)
      \;+\; \frac{\tau}{2}\,\Sigma(\theta)\,p'(\theta)
      \;=\; 0,
   \]
   so that
   \[
      \frac{p'(\theta)}{p(\theta)}
      \;=\; -\frac{2 L'(\theta)}{\tau\,\Sigma(\theta)}
            \;-\; \frac{\Sigma'(\theta)}{\Sigma(\theta)}
      \;=\; \frac{d}{d\theta}\Bigl[
                -\frac{2}{\tau}\,U(\theta) - \ln \Sigma(\theta)\Bigr].
   \]
   Integrating from a reference point $\theta_0$ and exponentiating gives
   \eqref{eq:stationary-1d-restate}; the rewriting
   $p_\tau \propto \exp(-(2/\tau)\Veff)$ follows by absorbing the $\ln
   \Sigma$ term into the exponent.  Normalizability of $Z_\tau$ is
   immediate from the exponential decay of $\exp(-(2/\tau)U)$ implied by
   \Cref{A:coerc} and \Cref{A:ellip}.
\end{proof}

\subsection{Symmetric two-well ratio formula}

\begin{lemma}[Symmetric two-well case]\label{lem:sym}
   Suppose, in addition to \Cref{A:two-wells}, that $L$ and $\Sigma$ are
   symmetric about $\theta_m := (\theta_A^* + \theta_B^*)/2$, in the sense
   that $L(\theta_m+u) = L(\theta_m-u)$ and $\Sigma(\theta_m+u)
   = \Sigma(\theta_m-u)$.  Then $U(\theta_A^*) = U(\theta_B^*)$, and
   \[
      \frac{p_\tau(\theta_A^*)}{p_\tau(\theta_B^*)}
      \;=\; \frac{\Sigma_B}{\Sigma_A}.
   \]
\end{lemma}

\begin{proof}
   The integrand $L'/\Sigma$ in $U$ is the ratio of an odd function ($L'$)
   and an even function ($\Sigma$) about $\theta_m$, hence is odd about
   $\theta_m$.  Choose the reference point $\theta_0=\theta_m$, so that
   $U(\theta_m) = 0$.  Then $U(\theta_m+u) = \int_0^u (L'/\Sigma)(\theta_m
   + s)\,ds$, which by the change of variables $s\mapsto -s$ and the
   oddness of the integrand equals $-\int_0^u (L'/\Sigma)(\theta_m -
   s)\,ds = -U(\theta_m - u)$.  Hence $U$ is odd about $\theta_m$:
   \begin{equation}\label{eq:U-odd}
      U(\theta_m+u) \;=\; -U(\theta_m-u).
   \end{equation}
   Setting $\theta_A^* = \theta_m - r$ and $\theta_B^* = \theta_m + r$
   (with $r>0$), \eqref{eq:U-odd} gives
   \begin{equation}\label{eq:U-A-B}
      U(\theta_A^*) \;=\; -U(\theta_B^*).
   \end{equation}
   Now we show in addition that both values equal zero.  Symmetry of $L$
   gives $L(\theta_A^*) = L(\theta_B^*)$, and since $L'(\theta_A^*) =
   L'(\theta_B^*) = 0$, the integral $U(\theta_B^*) - U(\theta_A^*) =
   \int_{\theta_A^*}^{\theta_B^*} L'(s)/\Sigma(s)\,ds$ can be split at the
   midpoint $\theta_m$:
   \[
      \int_{\theta_A^*}^{\theta_m} \frac{L'(s)}{\Sigma(s)}\,ds
      \;+\;
      \int_{\theta_m}^{\theta_B^*} \frac{L'(s)}{\Sigma(s)}\,ds
      \;=\; -U(\theta_A^*) + U(\theta_B^*)
      \;=\; 2\,U(\theta_B^*),
   \]
   using $U(\theta_m)=0$ and \eqref{eq:U-A-B}.  But by oddness of
   $L'/\Sigma$ about $\theta_m$, the two integrals on the left cancel,
   so $2U(\theta_B^*) = 0$, hence $U(\theta_A^*) = U(\theta_B^*) = 0$.

   Substituting into \eqref{eq:stationary-1d-restate}:
   \[
      \frac{p_\tau(\theta_A^*)}{p_\tau(\theta_B^*)}
      \;=\; \frac{\Sigma_B}{\Sigma_A}\,
            \exp\!\Bigl(-\tfrac{2}{\tau}\bigl(U(\theta_A^*)
                                              -U(\theta_B^*)\bigr)\Bigr)
      \;=\; \frac{\Sigma_B}{\Sigma_A}.
   \]
\end{proof}

\Cref{lem:exist-unique,lem:stationary-1d,lem:sym} together prove
\Cref{thm:closed-loop}~Part~I.  In the asymmetric case the formula
\eqref{eq:ratio-asym} is exactly \eqref{eq:stationary-1d-restate} evaluated
at the two wells; we record it as a corollary.

\begin{corollary}[General one-dimensional ratio]\label{cor:ratio-1d-gen}
   In the general $d=1$ setting under \Cref{A:C2}--\Cref{A:two-wells},
   \[
      \frac{p_\tau(\theta_A^*)}{p_\tau(\theta_B^*)}
      \;=\; \frac{\Sigma_B}{\Sigma_A}\,
            \exp\!\Bigl(-\tfrac{2}{\tau}\bigl(U(\theta_A^*)
                                                - U(\theta_B^*)\bigr)\Bigr).
   \]
\end{corollary}

%% =================================================================
\section{Proof of \Cref{thm:closed-loop}: Part~II (LOO gap)}
\label{sec:proof-2}
%% =================================================================

We now prove the leave-one-out generalization formula
\eqref{eq:LOO-gap}.  We work at a locally quadratic minimum $\theta^*$
with $H = H(\theta^*) > 0$ and $\Sigma = \Sigma(\theta^*)$.  Throughout
this section all $O(\cdot)$ constants depend only on the quantities listed
in the statement of Part~II.

\subsection{Perturbation of the minimizer}

\begin{lemma}[Leave-one-out shift]\label{lem:loo-shift}
   For each $j\in\{1,\dots,n\}$ there is a unique local minimizer
   $\theta^*_{-j}$ of $L_{-j}$ in a neighbourhood of $\theta^*$, and
   \begin{equation}\label{eq:loo-shift}
      \theta^*_{-j} - \theta^* \;=\;
      \frac{\nabla\ell_j(\theta^*)}{(n-1)\,H_{-j}}
      \;+\; O\!\Bigl(\frac{\norm{\nabla\ell_j(\theta^*)}^2}{n^2}\Bigr),
      \qquad H_{-j} \;:=\; \nabla^2 L_{-j}(\theta^*).
   \end{equation}
   Moreover $H_{-j} = H + O(1/n)$.
\end{lemma}

\begin{proof}
   Compute
   \[
      \nabla L_{-j}(\theta^*)
      \;=\; \frac{1}{n-1}\sum_{i\neq j} \nabla \ell_i(\theta^*)
      \;=\; \frac{n}{n-1}\nabla L(\theta^*)
            \;-\; \frac{1}{n-1}\nabla \ell_j(\theta^*)
      \;=\; -\,\frac{\nabla\ell_j(\theta^*)}{n-1},
   \]
   using $\nabla L(\theta^*) = 0$.  Similarly,
   \[
      \nabla^2 L_{-j}(\theta^*)
      \;=\; \frac{n}{n-1}\,H \;-\;
            \frac{\nabla^2 \ell_j(\theta^*)}{n-1}
      \;=\; H \;+\; \frac{1}{n-1}\bigl(H - \nabla^2 \ell_j(\theta^*)\bigr)
      \;=\; H + O(1/n).
   \]
   In a neighbourhood of $\theta^*$ where $\nabla^2 L_{-j}(\theta)\succ 0$
   (which holds for all sufficiently large $n$ by continuity), the
   implicit function theorem applied to $\nabla L_{-j}$ produces a unique
   solution $\theta^*_{-j}$ to $\nabla L_{-j}(\theta^*_{-j}) = 0$ with
   \[
      \theta^*_{-j} - \theta^*
      \;=\; -\,\bigl(\nabla^2 L_{-j}(\theta^*)\bigr)^{-1}\,
                  \nabla L_{-j}(\theta^*)
            \;+\; O\bigl(\norm{\nabla L_{-j}(\theta^*)}^2\bigr).
   \]
   Substituting the formulas above gives
   \eqref{eq:loo-shift}.
\end{proof}

\subsection{The pointwise generalization gap}

\begin{lemma}[Pointwise gap]\label{lem:pointwise-gap}
   With the notation of \Cref{lem:loo-shift},
   \begin{equation}\label{eq:pointwise-gap}
      \ell_j(\theta^*_{-j}) - \ell_j(\theta^*)
      \;=\; \frac{\norm{\nabla\ell_j(\theta^*)}^2}{(n-1)\,H_{-j}}
            \;+\; O\!\Bigl(\frac{\norm{\nabla\ell_j(\theta^*)}^3}{n^2}
                          \Bigr).
   \end{equation}
\end{lemma}

\begin{proof}
   Taylor-expand $\ell_j$ around $\theta^*$ to second order:
   \[
      \ell_j(\theta^*_{-j}) - \ell_j(\theta^*)
      \;=\; \nabla\ell_j(\theta^*)^{\!\top}(\theta^*_{-j}-\theta^*)
            \;+\; \tfrac{1}{2}(\theta^*_{-j}-\theta^*)^{\!\top}
                  \nabla^2\ell_j(\theta^*)
                  (\theta^*_{-j}-\theta^*)
            \;+\; O(\norm{\theta^*_{-j}-\theta^*}^3).
   \]
   Plugging in \eqref{eq:loo-shift} and using $\norm{\theta^*_{-j}
   -\theta^*} = O(\norm{\nabla\ell_j(\theta^*)}/n)$:
   \begin{align*}
      \nabla\ell_j(\theta^*)^{\!\top}(\theta^*_{-j}-\theta^*)
      &\;=\; \frac{\norm{\nabla\ell_j(\theta^*)}^2}{(n-1)\,H_{-j}}
             \;+\; O\!\Bigl(\frac{\norm{\nabla\ell_j(\theta^*)}^3}{n^2}\Bigr),\\
      \tfrac{1}{2}(\theta^*_{-j}-\theta^*)^{\!\top}
                  \nabla^2\ell_j(\theta^*)
                  (\theta^*_{-j}-\theta^*)
      &\;=\; O\!\Bigl(\frac{\norm{\nabla\ell_j(\theta^*)}^2}{n^2}\Bigr).
   \end{align*}
   Combining these and absorbing all lower-order remainders gives
   \eqref{eq:pointwise-gap}.
\end{proof}

\subsection{Averaging and the noise covariance}

\begin{proof}[Proof of \Cref{thm:closed-loop} Part~II]
   Average \eqref{eq:pointwise-gap} over $j$:
   \begin{align*}
      \Gap_{\LOO}(\theta^*)
      &\;=\; \frac{1}{n}\sum_{j=1}^n
             \frac{\norm{\nabla\ell_j(\theta^*)}^2}{(n-1)\,H_{-j}}
             \;+\; O\!\Bigl(\frac{1}{n^2}\Bigr)\\
      &\;=\; \frac{1}{(n-1)\,H}\,
             \cdot\,\frac{1}{n}\sum_{j=1}^n
                  \norm{\nabla\ell_j(\theta^*)}^2
             \;+\; O\!\Bigl(\frac{1}{n^2}\Bigr),
   \end{align*}
   where in the second line we replaced each $H_{-j}$ by $H$ at the cost
   of an $O(1/n)\cdot \E_j[\norm{\nabla\ell_j}^2]/(n-1)$ correction, which
   is $O(n^{-2})$ uniformly.  By \Cref{def:H-Sigma} and
   $\nabla L(\theta^*) = 0$,
   \[
      \frac{1}{n}\sum_{j=1}^n \norm{\nabla\ell_j(\theta^*)}^2
      \;=\; \tr\Bigl(\frac{1}{n}\sum_{j=1}^n
              \nabla\ell_j(\theta^*)\nabla\ell_j(\theta^*)^{\!\top}\Bigr)
      \;=\; \tr\Sigma(\theta^*).
   \]
   In dimension $d=1$ this is just $\Sigma(\theta^*)$, giving
   \eqref{eq:LOO-gap}.  In general dimension the identity becomes
   $\Gap_{\LOO}(\theta^*) = \tr(H^{-1}\Sigma)/(n-1) + O(n^{-2})$, where
   the trace formulation uses
   $\norm{\nabla\ell_j}^2_{H^{-1}} = \nabla\ell_j^{\!\top} H^{-1}
                                                          \nabla\ell_j$;
   this generalization is recorded in \Cref{prop:reparam}.
\end{proof}

\begin{remark}[Higher-order corrections]\label{rem:higher-order}
   The $O(n^{-2})$ remainder in \eqref{eq:LOO-gap} arises from the
   second-order Taylor remainder, the $H_{-j}\mapsto H$ replacement, and
   the $O(\norm{\nabla\ell_j}^2/n^2)$ correction in
   \Cref{lem:loo-shift}.  Each of these is bounded uniformly in $n$ by
   the data-dependent constants in the statement of Part~II.
\end{remark}

%% =================================================================
\section{Proof of \Cref{thm:closed-loop}: Part~III (implication)}
\label{sec:proof-3}
%% =================================================================

We now combine Parts~I and II.  Throughout this section we work with
the basin masses $M_\tau(A) = \int_{B_A} p_\tau$ and
$M_\tau(B) = \int_{B_B} p_\tau$ defined in Part~I, and we adopt the
equal-loss two-well setting $L(\theta_A^*) = L(\theta_B^*)$
(equivalently, $U(\theta_A^*) = U(\theta_B^*)$, see~\Cref{lem:sym} and
its derivation; the asymmetric case can be reduced to this one by
shifting $L$ by an additive constant within each basin).
Recall \Cref{A:sel}: $\Sigma_A/H_A < \Sigma_B/H_B$.

\subsection{Selection: SGD concentrates on the basin of $\theta_A^*$}

We apply Laplace's method to the integrals $M_\tau(A) = \int_{B_A}
p_\tau(\theta)\,d\theta$ and $M_\tau(B) = \int_{B_B} p_\tau(\theta)\,
d\theta$ directly.  Near the minimum $\theta_k^*$ ($k\in\{A,B\}$),
expand the effective potential $\Veff(\theta) = U(\theta) +
\tfrac{\tau}{2}\ln\Sigma(\theta)$ to second order: since $L'(\theta_k^*)
= 0$, the integrand $L'/\Sigma$ defining $U$ has $U'(\theta_k^*)=0$, and
$U''(\theta_k^*) = (L''/\Sigma)(\theta_k^*) = H_k/\Sigma_k$.  Therefore
\[
   p_\tau(\theta) \;=\; \frac{1}{Z_\tau\,\Sigma_k}
        \exp\!\Bigl(-\tfrac{2}{\tau}U(\theta_k^*)\Bigr)
        \exp\!\Bigl(-\tfrac{1}{\tau}\,\tfrac{H_k}{\Sigma_k}\,
                                 (\theta-\theta_k^*)^2\Bigr)
        (1 + O(\theta-\theta_k^*))
\]
on a neighbourhood of $\theta_k^*$.  Integrating against the Gaussian
factor over $B_k$ and using the standard Laplace estimate
$\int_{-\infty}^{\infty} e^{-\alpha u^2}\,du = \sqrt{\pi/\alpha}$ (the
contribution from outside any neighbourhood of $\theta_k^*$ is
exponentially smaller in $\tau$ by~\Cref{A:coerc} and~\Cref{A:ellip}),
\[
   M_\tau(k) \;=\; \frac{1}{Z_\tau\,\Sigma_k}\,
        \exp\!\Bigl(-\tfrac{2}{\tau}U(\theta_k^*)\Bigr)\,
        \sqrt{\pi\tau\,\Sigma_k/H_k}\,(1 + O(\sqrt\tau)).
\]
Taking the ratio at $A$ over $B$, the prefactor $1/Z_\tau$ cancels
and the exponentials combine.  Under the equal-loss assumption
$U(\theta_A^*) = U(\theta_B^*)$,
\begin{equation}\label{eq:basin-ratio-eqloss}
   \frac{M_\tau(A)}{M_\tau(B)}
   \;=\; \frac{\Sigma_B}{\Sigma_A}\,
         \sqrt{\frac{\Sigma_A/H_A}{\Sigma_B/H_B}}\,(1+O(\sqrt\tau))
   \;=\; \sqrt{\frac{\Sigma_B}{\Sigma_A}\,\frac{H_B}{H_A}}\,
                                                       (1+O(\sqrt\tau)).
\end{equation}
\emph{Symmetric two-well case ($H_A = H_B$, $\Sigma_A < \Sigma_B$).}
\Cref{A:sel} reduces to $\Sigma_A < \Sigma_B$, and
\eqref{eq:basin-ratio-eqloss} simplifies to $M_\tau(A)/M_\tau(B) =
\sqrt{\Sigma_B/\Sigma_A}\,(1+O(\sqrt\tau)) > 1$ for $\tau$ small.

\emph{General equal-loss case.}  In general,
\eqref{eq:basin-ratio-eqloss} shows that $M_\tau(A)/M_\tau(B) > 1$ iff
$\Sigma_B H_B > \Sigma_A H_A$.  This condition is \emph{not} implied by
\Cref{A:sel} alone (only $\Sigma/H$ smaller at $A$); accordingly we
adopt for the implication (i) the explicit
\begin{equation}\label{eq:sigmaH-product}
   \Sigma_B\,H_B \;>\; \Sigma_A\,H_A,
\end{equation}
which holds automatically in the symmetric case ($H_A=H_B$) and
otherwise must be checked.  Under \eqref{eq:sigmaH-product}, for
$\tau$ sufficiently small, $M_\tau(A)/M_\tau(B) > 1$, proving
conclusion~(i).  We note that the predictive content of Part~III is
unaffected: parts (ii) and (iii) below depend only on \Cref{A:sel}, and
the only role of (i) is to confirm that SGD's stationary mass is
concentrated on the better basin in the symmetric case to which
empirical tests of the theorem typically apply.

\subsection{Generalization: $\theta_A^*$ has smaller LOO gap}

Apply Part~II at both minima:
\begin{align*}
   \Gap_{\LOO}(\theta_A^*) &\;=\;
       \frac{\Sigma_A}{(n-1)\,H_A}\;+\; O(n^{-2}),\\
   \Gap_{\LOO}(\theta_B^*) &\;=\;
       \frac{\Sigma_B}{(n-1)\,H_B}\;+\; O(n^{-2}).
\end{align*}
Subtracting,
\[
   \Gap_{\LOO}(\theta_B^*) - \Gap_{\LOO}(\theta_A^*)
   \;=\; \frac{1}{n-1}\Bigl(\tfrac{\Sigma_B}{H_B}
                                       - \tfrac{\Sigma_A}{H_A}\Bigr)
            \;+\; O(n^{-2}).
\]
By \Cref{A:sel}, the leading term is strictly positive.  This proves
(ii) of Part~III, with explicit margin
$\bigl(\Sigma_B/H_B - \Sigma_A/H_A\bigr)/(n-1)$.

\subsection{Test loss bound}

By the standard $\E[\Gap_{\LOO}] = \E[L_{\text{pop}}(\theta^*) -
L(\theta^*)]$ identity \cite{BousquetElisseeff2002}, the LOO gap is an
unbiased estimator of the population--training gap, provided that the
data--model interaction does not introduce a systematic reversal between
LOO behaviour and population behaviour.  We make this explicit by
invoking \Cref{A:F4-neg}: under the negation of failure mode~F4, the
sign of the population-gap difference between two minima agrees with
the sign of the LOO-gap difference, up to $o(1/n)$.  Combined with
$L(\theta_A^*) = L(\theta_B^*)$ (equal-loss two wells), the test loss
bound at $\theta_A^*$ minus that at $\theta_B^*$ equals
\eqref{eq:gap-margin} up to $O(n^{-2})$.  This proves (iii) of
Part~III.  Without \Cref{A:F4-neg}, parts (i) and (ii) still hold, but
(iii) need not: this is the precise location at which the closed loop
of \Cref{thm:closed-loop} can be broken by F4, as discussed in
\Cref{sec:discussion}.
\qed

\begin{remark}[Why this is a closed loop]\label{rem:closed-loop}
   The implication is genuinely a closed loop: Part~I selects the minimum
   that has lower $\Sigma$ (in the symmetric case, lower
   $\Sigma_A/H_A$ in the asymmetric case via Laplace), Part~II shows
   that lower $\Sigma/H$ corresponds to a lower LOO gap, and Part~III
   shows that the same intrinsic quantity drives both.  The chain
   \[
      \mathrm{SGD\ stationary\ measure}
      \;\longrightarrow\; \text{low } \Sigma/H
      \;\longrightarrow\; \text{low LOO gap}
      \;\longrightarrow\; \text{low test loss}
   \]
   is closed by \Cref{prop:reparam} (which guarantees that the
   intermediate quantity is intrinsic).  The chain breaks if and only if
   one of the failure modes F1--F4 of \cite{paperB} is active; we make
   this precise in \Cref{sec:discussion}.
\end{remark}

%% =================================================================
\section{Multivariate extension via Freidlin--Wentzell}\label{sec:multi}
%% =================================================================

In dimension $d>1$ the FPE \eqref{eq:fpe} need not satisfy detailed balance,
and the closed-form Gibbs density \eqref{eq:stationary-1d} need not exist.
The correct generalization is to replace $U(\theta)$ by the
\emph{Freidlin--Wentzell quasipotential} \cite{FreidlinWentzell2012}.

\subsection{Quasipotential definition and HJ equation}

\begin{definition}[Quasipotential]\label{def:quasi}
   Fix a stable equilibrium $\theta_0\in\Theta$ of $-\nabla L$.  The
   Freidlin--Wentzell quasipotential at $\theta\in\Theta$ is
   \begin{equation}\label{eq:quasipotential}
      V(\theta) \;:=\;
      \inf\Bigl\{\tfrac{1}{2}\!\int_0^T
            (\dot\phi + \nabla L(\phi))^{\!\top}\,
            \Sigma(\phi)^{-1}\,(\dot\phi + \nabla L(\phi))\,dt
         \;:\;
         T>0,\;\phi(0)=\theta_0,\;\phi(T)=\theta\Bigr\}.
   \end{equation}
\end{definition}

\begin{lemma}[Hamilton--Jacobi characterization]\label{lem:HJ}
   Under \Cref{A:C2}--\Cref{A:ellip}, the function $V$ in
   \Cref{def:quasi} is locally Lipschitz on $\Theta$, and is the unique
   viscosity solution (in the Lipschitz sense) of the stationary
   Hamilton--Jacobi equation
   \begin{equation}\label{eq:HJ}
      -\nabla L(\theta)^{\!\top}\,\nabla V(\theta)
      \;+\; \tfrac{1}{2}\,\nabla V(\theta)^{\!\top}\,\Sigma(\theta)\,
            \nabla V(\theta)
      \;=\; 0,
      \qquad V(\theta_0) = 0,\;\; V\geq 0,
   \end{equation}
   with the prescribed boundary value at $\theta_0$.  Moreover $V$ is
   $C^1$ in a neighbourhood of $\theta_0$ (where $V$ admits a smooth
   quadratic expansion), but globally only Lipschitz regularity is
   guaranteed---$V$ may fail to be $C^1$ along the cut locus of the
   underlying minimum-action geodesic flow.
\end{lemma}

\begin{proof}
   Standard Freidlin--Wentzell theory; see
   \cite[Chapter~4]{FreidlinWentzell2012} or
   \cite[Theorem~5.7.11]{DemboZeitouni2010}.  Local Lipschitz regularity
   of $V$ is a consequence of the uniform ellipticity of $\Sigma$ and
   the smoothness of $\nabla L$.  Local $C^1$ regularity near
   $\theta_0$ follows from the implicit function theorem applied to
   the Hamiltonian flow associated with \eqref{eq:HJ} at the stable
   equilibrium $(\theta_0,0)$.  Global $C^1$ regularity is \emph{not}
   guaranteed by Freidlin--Wentzell theory; it requires additional
   structural conditions on $(L,\Sigma)$ that are not used in our
   theorem.
\end{proof}

\begin{proposition}[Stationary measure with quasipotential]
   \label{prop:stationary-multi}
   Under \Cref{A:C2}--\Cref{A:ellip}, the FPE \eqref{eq:fpe} has a unique
   stationary density $p_\tau$, and as $\tau\to 0$,
   \[
      -\,\tau \ln p_\tau(\theta) \;\longrightarrow\; 2\,V(\theta)
      \quad \text{uniformly on compact subsets of }\Theta,
   \]
   where $V$ is the quasipotential of \Cref{def:quasi}.  The leading-order
   ratio between two stable equilibria is
   \[
      \frac{p_\tau(\theta_A^*)}{p_\tau(\theta_B^*)}
      \;=\; \exp\!\Bigl(-\tfrac{2}{\tau}\bigl(V(\theta_A^*)
                                              - V(\theta_B^*)\bigr)\Bigr)\,
            (1 + o(1)).
   \]
\end{proposition}

\begin{proof}
   This is the Freidlin--Wentzell large-deviations principle for invariant
   measures of small-noise diffusions; see
   \cite[Theorem~4.4.1]{FreidlinWentzell2012} and
   \cite[Theorem~5.7.12]{DemboZeitouni2010}.  Existence and uniqueness
   of the stationary density follow as in \Cref{lem:exist-unique}.
\end{proof}

\subsection{A 2D worked example}\label{sec:2d-example}

We now state and prove \Cref{thm:multi}, in which detailed balance fails
but the quasipotential admits a closed-form expansion to fourth order.

\begin{theorem}[Multivariate extension; 2D worked example]\label{thm:multi}
   Let $\Theta = \R^2$, $L(\theta_1,\theta_2) = \tfrac{1}{2}(\theta_1^2
   + \theta_2^2)$, and
   \[
      \Sigma(\theta) \;=\; \diag\bigl(1 + \theta_2^2,\,1\bigr).
   \]
   Then:
   \begin{enumerate}[label=\textnormal{(\alph*)}]
   \item Detailed balance fails: the vector field
       $b(\theta) := \Sigma(\theta)^{-1}\nabla L(\theta)$ has nonzero
       curl, namely
       $\partial_{\theta_2} b_1 - \partial_{\theta_1} b_2
                 = -2\theta_1\theta_2/(1+\theta_2^2)^2 \neq 0$.
   \item The quasipotential $V$ admits a Hamilton--Jacobi expansion
       \begin{equation}\label{eq:V-expand}
          V(\theta_1,\theta_2)
          \;=\; \theta_1^2 + \theta_2^2
              \;-\; \tfrac{1}{2}\,\theta_1^2 \theta_2^2
              \;+\; O(\norm\theta^6),
       \end{equation}
       which is the unique even $C^1$ solution of \eqref{eq:HJ} with
       $V(0)=0$ and $V\geq 0$ to the displayed order.
   \item The leading correction $-\tfrac{1}{2}\theta_1^2\theta_2^2$ is
       a saddle perturbation: $V$ is reduced along the diagonal where
       both $|\theta_1|$ and $|\theta_2|$ are large simultaneously.
       Consequently the stationary measure places more mass along this
       diagonal direction than the naive Gaussian $\exp(-2L/\tau)$ would
       predict.
   \end{enumerate}
\end{theorem}

\begin{proof}
   \emph{(a)}  Compute
   \[
      b(\theta) \;=\; \Sigma(\theta)^{-1}\,(\theta_1,\theta_2)
                  \;=\; \Bigl(\tfrac{\theta_1}{1+\theta_2^2},\,\theta_2\Bigr).
   \]
   Then $\partial_{\theta_2} b_1 = -2\theta_1\theta_2/(1+\theta_2^2)^2$
   and $\partial_{\theta_1} b_2 = 0$.  Hence
   $\curl b = -2\theta_1\theta_2/(1+\theta_2^2)^2 \neq 0$ on the open set
   $\{\theta_1\theta_2\neq 0\}$, so detailed balance fails.

   \emph{(b)} The HJ equation \eqref{eq:HJ} reads
   \begin{equation}\label{eq:HJ-2d}
      -(\theta_1\partial_{\theta_1}V + \theta_2\partial_{\theta_2} V)
      \;+\; \tfrac{1}{2}\bigl((1+\theta_2^2)(\partial_{\theta_1}V)^2
                              + (\partial_{\theta_2}V)^2\bigr)
      \;=\; 0.
   \end{equation}
   We seek an even formal power-series solution
   $V = \sum_{k\geq 1} V_{2k}$, where $V_{2k}$ is homogeneous of degree
   $2k$.  At leading order $V_2$ satisfies
   \[
      -2 V_2 \;+\; \tfrac{1}{2}\bigl(\norm{\nabla V_2}^2\bigr) \;=\; 0
      \;\;\Rightarrow\;\; V_2(\theta) = \theta_1^2 + \theta_2^2,
   \]
   the unique even quadratic such that $V_2(0)=0$, $V_2\geq 0$.  At order
   $4$, write $V_4(\theta) = a\theta_1^4 + b\theta_1^2\theta_2^2
   + c\theta_2^4$.  Substituting $V = V_2 + V_4$ into \eqref{eq:HJ-2d}
   and collecting degree-$4$ terms requires three contributions:
   \begin{enumerate}[label=(\roman*),leftmargin=2.4em]
   \item the drift part
       $-(\theta_1\partial_1 V_4 + \theta_2\partial_2 V_4)
         = -(4a\theta_1^4 + 4b\theta_1^2\theta_2^2 + 4c\theta_2^4)$,
   \item the diffusion cross term
       $\nabla V_2\cdot\nabla V_4
         = 2\theta_1(4a\theta_1^3 + 2b\theta_1\theta_2^2)
         + 2\theta_2(2b\theta_1^2\theta_2 + 4c\theta_2^3)
         = 8a\theta_1^4 + 8b\theta_1^2\theta_2^2 + 8c\theta_2^4$,
   \item the $\Sigma$-correction at order $4$ from the
       $\theta_2^2(\partial_1 V)^2$ term in \eqref{eq:HJ-2d}, which at
       leading order contributes
       $\tfrac{1}{2}\theta_2^2(\partial_1 V_2)^2
        = \tfrac{1}{2}\theta_2^2\cdot 4\theta_1^2 = 2\theta_1^2\theta_2^2$.
   \end{enumerate}
   (No order-$4$ contribution arises from $(\partial_2 V_4)^2$ or
   $(\partial_1 V_4)^2$, since these are of order $6$.)  Summing the
   three contributions, \eqref{eq:HJ-2d} at order four becomes
   \[
      -(4a\theta_1^4 + 4b\theta_1^2\theta_2^2 + 4c\theta_2^4)
      \;+\; \bigl(8a\theta_1^4 + 8b\theta_1^2\theta_2^2 + 8c\theta_2^4\bigr)/2\cdot 2
      \;+\; 2\theta_1^2\theta_2^2 \;=\; 0,
   \]
   which simplifies (since $\tfrac{1}{2}\cdot(8a,8b,8c) = (4a,4b,4c)$) to
   \begin{equation}\label{eq:order-4}
      4a\,\theta_1^4 \;+\; (4b+2)\,\theta_1^2\theta_2^2
                                       \;+\; 4c\,\theta_2^4 \;=\; 0.
   \end{equation}
   Matching coefficients of the three independent monomials gives
   $a = 0$, $b = -\tfrac{1}{2}$, $c = 0$, hence
   \[
      V_4(\theta) \;=\; -\tfrac{1}{2}\,\theta_1^2\theta_2^2.
   \]
   This proves \eqref{eq:V-expand}: $V = \theta_1^2 + \theta_2^2
   - \tfrac{1}{2}\theta_1^2\theta_2^2 + O(\norm\theta^6)$.

   We now verify directly that \eqref{eq:V-expand} solves
   \eqref{eq:HJ-2d} to order four.  With $V = \theta_1^2 + \theta_2^2
   - \tfrac{1}{2}\theta_1^2\theta_2^2$, $\partial_1 V = 2\theta_1
   - \theta_1\theta_2^2$ and $\partial_2 V = 2\theta_2 -
   \theta_1^2\theta_2$.  Then
   \begin{align*}
      -\theta_1\partial_1 V - \theta_2\partial_2 V
      &= -\theta_1(2\theta_1 - \theta_1\theta_2^2)
                  -\theta_2(2\theta_2 - \theta_1^2\theta_2)\\
      &= -2(\theta_1^2+\theta_2^2) + 2\theta_1^2\theta_2^2,\\
      \tfrac{1}{2}(1+\theta_2^2)(\partial_1 V)^2
      &= \tfrac{1}{2}(1+\theta_2^2)\bigl(4\theta_1^2
                              - 4\theta_1^2\theta_2^2
                              + \theta_1^2\theta_2^4\bigr)\\
      &= 2\theta_1^2 + 2\theta_1^2\theta_2^2 - 2\theta_1^2\theta_2^2
                                            + O(\norm\theta^6)\\
      &= 2\theta_1^2 + O(\norm\theta^6),\\
      \tfrac{1}{2}(\partial_2 V)^2
      &= \tfrac{1}{2}(4\theta_2^2 - 4\theta_1^2\theta_2^2
                                  + \theta_1^4\theta_2^2)
       = 2\theta_2^2 - 2\theta_1^2\theta_2^2 + O(\norm\theta^6).
   \end{align*}
   Summing,
   \[
      -2\theta_1^2 - 2\theta_2^2 + 2\theta_1^2\theta_2^2
      + 2\theta_1^2 + 2\theta_2^2 - 2\theta_1^2\theta_2^2
      + O(\norm\theta^6) \;=\; O(\norm\theta^6),
   \]
   confirming \eqref{eq:HJ-2d} to order four.  Uniqueness among even
   $C^1$ solutions of polynomial form with $V(0)=0$, $V\geq 0$ follows
   from the fact that \eqref{eq:order-4} uniquely determines the three
   coefficients $(a,b,c)$, and analogously at every higher even order
   the matching conditions are linear with unique solution.

   \emph{(c)} The Gaussian density $\exp(-2L/\tau) = \exp(-(\theta_1^2 +
   \theta_2^2)/\tau)$ corresponds to $V_{\text{naive}} =
   \theta_1^2+\theta_2^2$.  The correction $-\tfrac{1}{2}\theta_1^2
   \theta_2^2$ in \eqref{eq:V-expand} is non-positive, so $V \leq
   V_{\text{naive}}$, with equality on the axes $\{\theta_1=0\}$ and
   $\{\theta_2=0\}$ and strict inequality in the off-axis region.
   Consequently $\exp(-2V/\tau) \geq \exp(-2 V_{\text{naive}}/\tau)$,
   establishing (c).
\end{proof}

\begin{remark}[Curl drives the correction]\label{rem:curl-drives}
   The order-four correction $-\tfrac{1}{2}\theta_1^2\theta_2^2$ in
   \eqref{eq:V-expand} arises from the same term that produces nonzero
   curl in part~(a).  More precisely, the symmetric part of
   $\Sigma^{-1}\nabla\nabla L$ contributes to $V_2$ while the
   antisymmetric part---namely $\curl b$---is the source term for $V_4$
   through the cross-coupling in \eqref{eq:HJ-2d}.  In dimension $d=1$
   curl vanishes identically, recovering the closed-form
   \eqref{eq:stationary-1d}.
\end{remark}

%% =================================================================
\section{Reparameterization invariance}\label{sec:reparam}
%% =================================================================

\begin{proposition}[Reparameterization invariance]\label{prop:reparam}
   Let $\Phi:\Theta'\to \Theta$ be a $C^2$ diffeomorphism, write
   $\theta = \Phi(\phi)$, and let $J_\Phi(\phi)$ be the Jacobian.  Define
   the pulled-back loss $\widetilde L := L\circ\Phi$ and the pulled-back
   per-sample losses $\widetilde\ell_i := \ell_i\circ\Phi$.  At a critical
   point $\phi^* := \Phi^{-1}(\theta^*)$ of $\widetilde L$:
   \begin{enumerate}[label=\textnormal{(\alph*)}]
   \item The Hessian transforms as
       $\widetilde H(\phi^*) = J_\Phi(\phi^*)^{\!\top}\,H(\theta^*)\,
                                                            J_\Phi(\phi^*)$.
   \item The noise covariance transforms as
       $\widetilde\Sigma(\phi^*) = J_\Phi(\phi^*)^{\!\top}\,\Sigma(\theta^*)\,
                                                            J_\Phi(\phi^*)$.
   \item Consequently $\tr(\widetilde H^{-1}\widetilde\Sigma)
                        = \tr(H^{-1}\Sigma)$, and in particular in
       dimension $d=1$ the scalar ratio $\Sigma/H$ is invariant.
   \end{enumerate}
\end{proposition}

\begin{proof}
   Write $J=J_\Phi(\phi^*)$.  At any $\phi$,
   $\widetilde\ell_i(\phi) = \ell_i(\Phi(\phi))$, so by the chain rule
   \[
      \nabla_\phi \widetilde\ell_i(\phi)
      \;=\; J_\Phi(\phi)^{\!\top}\,\nabla_\theta \ell_i(\Phi(\phi)).
   \]
   At $\phi^*$ this becomes
   $\nabla_\phi \widetilde\ell_i(\phi^*)
   = J^{\!\top}\,\nabla_\theta\ell_i(\theta^*)$.
   Therefore
   \[
      \widetilde\Sigma(\phi^*)
      \;=\; \frac{1}{n}\sum_{i=1}^n \nabla_\phi\widetilde\ell_i(\phi^*)
                                   \nabla_\phi\widetilde\ell_i(\phi^*)^{\!\top}
      \;=\; J^{\!\top}\Bigl(\frac{1}{n}\sum_{i=1}^n
                              \nabla_\theta\ell_i\nabla_\theta\ell_i^{\!\top}
                            \Bigr) J
      \;=\; J^{\!\top}\Sigma(\theta^*) J,
   \]
   using $\nabla_\theta L(\theta^*)=0$ to drop the cross term.  This
   proves~(b).

   For~(a), differentiate $\nabla_\phi \widetilde L(\phi)
   = J_\Phi(\phi)^{\!\top}\nabla_\theta L(\Phi(\phi))$ once more with
   respect to $\phi$:
   \[
      \nabla_\phi^2 \widetilde L(\phi)
      \;=\; J_\Phi(\phi)^{\!\top}\,\nabla_\theta^2 L(\Phi(\phi))\,J_\Phi(\phi)
            \;+\; \sum_{a=1}^d
              \bigl[\nabla_\theta L(\Phi(\phi))\bigr]_a\,
              \nabla_\phi^2 \Phi_a(\phi).
   \]
   At $\phi^*$, $\nabla_\theta L(\theta^*) = 0$, so the second sum
   vanishes and we obtain
   $\widetilde H(\phi^*) = J^{\!\top} H(\theta^*) J$, proving (a).

   For~(c), with $\widetilde H = J^{\!\top} H J$ and
   $\widetilde\Sigma = J^{\!\top}\Sigma J$, and assuming $H,J$ invertible
   (so $\widetilde H$ is also invertible):
   \[
      \widetilde H^{-1} \widetilde \Sigma
      \;=\; (J^{\!\top}HJ)^{-1}(J^{\!\top}\Sigma J)
      \;=\; J^{-1} H^{-1} J^{-\!\top}\,J^{\!\top}\Sigma J
      \;=\; J^{-1} H^{-1}\Sigma\, J.
   \]
   Hence
   \[
      \tr(\widetilde H^{-1}\widetilde\Sigma)
      \;=\; \tr(J^{-1} H^{-1}\Sigma\, J)
      \;=\; \tr(H^{-1}\Sigma\, J\, J^{-1})
      \;=\; \tr(H^{-1}\Sigma),
   \]
   using cyclicity of the trace.  In dimension $d=1$, $J$ and $H$ are
   nonzero scalars and $\widetilde\Sigma/\widetilde H
   = (J^2 \Sigma)/(J^2 H) = \Sigma/H$.
\end{proof}

\begin{remark}[The LOO gap is intrinsic]\label{rem:loo-intrinsic}
   Combined with \Cref{thm:closed-loop} Part~II, \Cref{prop:reparam} (c)
   implies that, in dimension $d>1$, the LOO gap formula
   \eqref{eq:LOO-gap} reads
   $\Gap_{\LOO}(\theta^*) = \tr(H^{-1}\Sigma)/(n-1) + O(n^{-2})$, and is
   invariant under reparameterization.  In contrast, the curvature
   alone (e.g.\ a single eigenvalue or $\det H$) is not invariant
   \cite{Dinh2017}, and so cannot serve, by itself, as an intrinsic
   complexity measure.  The empirical phenomenon that large-batch SGD
   tends to find ``sharp'' minima with worse generalization
   \cite{KeskarMudigereSmelyanskiyTangNocedal2017} is consistent with
   the present picture once curvature is replaced by the
   reparameterization-invariant quantity $\tr(H^{-1}\Sigma)$.
\end{remark}

%% =================================================================
\section{Discussion}\label{sec:discussion}
%% =================================================================

\subsection{Boundary with paper~B's failure modes}

The closed-loop \Cref{thm:closed-loop} fails exactly when one of the four
failure modes catalogued in \cite{paperB} is active, in the following
precise correspondence.

\paragraph{F1 (insufficient mixing).}
The proof of Part~I uses the FPE stationary density as a proxy for the
SGD long-time average; this requires the mixing time $\tau_{\mathrm{mix}}$
between the two basins to be much smaller than the runtime $T$.
When $\tau_{\mathrm{mix}} \gg T$ (Kramers regime,
$\tau_{\mathrm{mix}}\sim \exp(2\Delta V/\tau)$ for barrier height
$\Delta V$), SGD does not reach the stationary measure and Part~I's
prediction does not apply.

\paragraph{F2 (discrete instability).}
The SDE approximation \eqref{eq:sgd-sde} is valid only for $\eta$ below
the discrete stability threshold $2/\lambda_{\max}(H)$ \cite{LiTaiE2017};
above it the SGD recursion diverges and no stationary measure exists.

\paragraph{F3 (degenerate selection).}
When $\Sigma_A/H_A = \Sigma_B/H_B$, \Cref{A:sel} fails: the implication of
Part~III is vacuous and the test-loss bounds at the two minima coincide
to leading order.

\paragraph{F4 (noise--generalization reversal).}
The implication \emph{(iii)} of Part~III rests on the alignment of low
$\Sigma/H$ with low population loss.  Paper~B exhibits explicit examples
in which this alignment is reversed: a memorizing minimum may have
$\Sigma_B \ll \Sigma_A$ at a generalizing minimum because all per-sample
gradients agree at the memorized solution.  When F4 is active,
\Cref{thm:closed-loop} Parts~I and II still hold individually, but their
composition does not yield a generalization guarantee.  In this sense F4
is the one fundamental obstruction; F1--F3 can be removed by tuning of
$\eta,B,T$, but F4 is a property of the data--architecture pair.

\subsection{Relation to paper~A (interpolation regime)}

Paper~A \cite{paperA} studies the interpolation regime where
$\Sigma(\theta^*)=0$ at every interpolating solution.  There the
Fokker--Planck argument of \Cref{sec:proof-1} is vacuous (the diffusion
degenerates), and a discrete Lyapunov-exponent analysis takes its place.
The quantities that play the role of $\Sigma/H$ are:
\begin{itemize}
\item In paper~A, the per-sample curvature $h_{\max} = \max_i
   \norm{\nabla_\theta f_{\theta^*}(x_i)}^2$ and the Jacobian
   complexity $\mathcal{C} = \norm{J}_F/\sqrt n$.  Stability requires
   $\eta\,h_{\max} < 2$, and stable solutions have $\mathcal{C}\leq
   \sqrt{2/\eta}$.
\item In the present paper~C, the noise-to-curvature ratio $\Sigma/H$
   (or $\tr(H^{-1}\Sigma)$).
\end{itemize}
Both are intrinsic (\Cref{prop:reparam} for the present paper;
\cite[Sec.~6]{paperA} for the discrete case), and both deliver
$O(\rho/\sqrt{\eta n})$- or $O(1/n)$-type generalization bounds.  The
unifying picture is that SGD's noise selects against an intrinsic
complexity measure---multiplicative noise covariance in the
non-interpolation regime, per-sample Jacobian curvature in the
interpolation regime---and that this selection is reflected in
sample-based generalization (LOO in the present setting; Rademacher in
\cite{paperA}).  The discrete-stability viewpoint in
\cite{WuSu2023} provides an alternative (dynamical-stability-based)
treatment of implicit regularization in the same regime addressed by
paper~A; our paper~C complements this by treating the
non-interpolation regime via Fokker--Planck.

\subsection{Open questions}

\begin{enumerate}[leftmargin=2em]
\item \emph{Higher-order quasipotential.}  \Cref{thm:multi} computes the
   quasipotential to order $\norm\theta^4$.  Beyond order $4$, the
   HJ equation \eqref{eq:HJ-2d} becomes substantially more complex; a
   global structural understanding of $V$ in $d=2$ is open.
\item \emph{Multivariate detailed balance.}  Identifying the largest
   class of $(L,\Sigma)$ for which detailed balance holds beyond the
   scalar-isotropic case is open.  \cite{ChaudhariSoatto2018} treats the
   constant-$\Sigma$ case; the state-dependent case requires curl
   conditions of the form $\curl(\Sigma^{-1}\nabla L)=0$, whose general
   structure is unknown.
\item \emph{Sharp test-loss bound.}  The implication (iii) of Part~III
   gives a relative bound but not an absolute one.  Combining
   \Cref{thm:closed-loop} with concentration inequalities to obtain an
   absolute population-loss bound (in expectation or with high probability)
   would close a remaining gap.
\item \emph{F4 detection.}  When does F4 occur in deep networks?  Paper~B
   \cite{paperB} provides constructions but the empirical prevalence is
   open.
\end{enumerate}

\subsection{Testable predictions}

Paper-C's theorem suggests the following testable predictions, which
can be empirically verified with standard SGD experiments:
\begin{enumerate}[leftmargin=2em]
\item For two locally quadratic minima with measurable $H$ and
   $\Sigma$, the ratio of stationary basin masses under SGD should
   match $\Sigma_B/\Sigma_A$ in the symmetric-well case ($H_A=H_B$),
   per the basin-mass formula \eqref{eq:basin-ratio-eqloss}; in the
   asymmetric equal-loss case the prediction generalizes to
   $\sqrt{(\Sigma_B/\Sigma_A)(H_B/H_A)}$ via Laplace approximation.
\item The LOO gap, which is computable in $O(n)$ time at any local
   minimum, should match $\Sigma/(H(n-1))$ to leading order.  If it does
   not, F4 is a possible explanation.
\item Reparameterizing the parameter space (e.g.\ rescaling layers in a
   neural network) should leave $\tr(H^{-1}\Sigma)$ unchanged at any
   minimum, even though both $H$ and $\Sigma$ individually rescale.
\end{enumerate}

%% =================================================================
%% REFERENCES
%% =================================================================
\begin{thebibliography}{99}

\bibitem{paperA}
L.~Chang.
\emph{Discrete Lyapunov stability of SGD in the interpolation regime:
   per-sample analysis and generalization bounds}.
Companion paper, 2026.

\bibitem{paperB}
L.~Chang.
\emph{When does SGD fail to find good solutions?  Four failure modes and
   necessary conditions for implicit regularization}.
Companion paper, 2026.

\bibitem{LiTaiE2017}
Q.~Li, C.~Tai, and W.~E.
\emph{Stochastic modified equations and adaptive stochastic gradient
   algorithms}.
In Proceedings of ICML, pp.\ 2101--2110, 2017.

\bibitem{MandtHoffmanBlei2017}
S.~Mandt, M.~D.~Hoffman, and D.~M.~Blei.
\emph{Stochastic gradient descent as approximate Bayesian inference}.
Journal of Machine Learning Research \textbf{18}(134):1--35, 2017.

\bibitem{SmithLe2018}
S.~L.~Smith and Q.~V.~Le.
\emph{A Bayesian perspective on generalization and stochastic gradient
   descent}.
In Proceedings of ICLR, 2018.

\bibitem{ChaudhariSoatto2018}
P.~Chaudhari and S.~Soatto.
\emph{Stochastic gradient descent performs variational inference,
   converges to limit cycles for deep networks}.
In Proceedings of ICLR, 2018.

\bibitem{HaoChen2021}
J.~Z.~HaoChen, C.~Wei, J.~D.~Lee, and T.~Ma.
\emph{Shape matters: understanding the implicit bias of the noise
   covariance}.
In Proceedings of COLT, pp.\ 2315--2357, 2021.

\bibitem{FreidlinWentzell2012}
M.~I.~Freidlin and A.~D.~Wentzell.
\emph{Random Perturbations of Dynamical Systems}, third edition.
Grundlehren der mathematischen Wissenschaften, vol.\ 260,
Springer, 2012.

\bibitem{DemboZeitouni2010}
A.~Dembo and O.~Zeitouni.
\emph{Large Deviations Techniques and Applications}, second edition.
Stochastic Modelling and Applied Probability, vol.\ 38, Springer, 2010.

\bibitem{Pavliotis2014}
G.~A.~Pavliotis.
\emph{Stochastic Processes and Applications: Diffusion Processes, the
   Fokker--Planck and Langevin Equations}.
Texts in Applied Mathematics, vol.\ 60, Springer, 2014.

\bibitem{BousquetElisseeff2002}
O.~Bousquet and A.~Elisseeff.
\emph{Stability and generalization}.
Journal of Machine Learning Research \textbf{2}:499--526, 2002.

\bibitem{MohriRostamizadehTalwalkar2018}
M.~Mohri, A.~Rostamizadeh, and A.~Talwalkar.
\emph{Foundations of Machine Learning}, second edition.
MIT Press, 2018.

\bibitem{Dinh2017}
L.~Dinh, R.~Pascanu, S.~Bengio, and Y.~Bengio.
\emph{Sharp minima can generalize for deep nets}.
In Proceedings of ICML, pp.\ 1019--1028, 2017.

\bibitem{BottouCurtisNocedal2018}
L.~Bottou, F.~E.~Curtis, and J.~Nocedal.
\emph{Optimization methods for large-scale machine learning}.
SIAM Review \textbf{60}(2):223--311, 2018.

\bibitem{WuSu2023}
L.~Wu and W.~J.~Su.
\emph{The implicit regularization of dynamical stability in stochastic
   gradient descent}.
In Proceedings of ICML, 2023.

\bibitem{KeskarMudigereSmelyanskiyTangNocedal2017}
N.~S.~Keskar, D.~Mudigere, J.~Nocedal, M.~Smelyanskiy, and P.~T.~P.~Tang.
\emph{On large-batch training for deep learning: generalization gap and
   sharp minima}.
In Proceedings of ICLR, 2017.

\end{thebibliography}

\end{document}
